%!TEX root = ../chapter06.tex 
%----------------------------------------------------------------------------------------
%	
%----------------------------------------------------------------------------------------
%----------------------------------------------------------------------------------------
%	 
%----------------------------------------------------------------------------------------


\section*{Vers la géométrie repérée}
\begin{figure*}[hb]\centering
	\begin{tikzpicture}[scale=1]
		\tkzSetUpPoint[size=3, fill=black]
		\tkzfctset{tan style/.style={->,>=latex, color=red, line width=1.2pt}}   
		\tkzSetUpLine[line width= 1.2pt, add=.2 and .2]  
		\begin{scope}[dashed] 
			\tkzGrid[ xstep=1, ystep=1](-6,-5)(12,7)
		\end{scope}   
		%	
		%	
		\tkzInit[xmin=-6,xmax=12.0,ymin=-5.0,ymax=7,xstep=1.0, ystep=1.0]; 
		\tkzDrawX[ font=\small, line width=2pt, label=$x$ , above =5pt ] 
		\tkzDrawY[ font=\small, line width=2pt , label={$y$} , right =5pt ]  
		\tkzLabelXY[orig = false,  font=\scriptsize] 
		\tkzRep[ xlabel=,ylabel=, font=\scriptsize, line width = 2pt,xnorm=1, ynorm=1]  
		\tkzDefPoints{-5/-3/A,1/5/B, -1/-4/C,4/-4/D,6/-1/E,11/-1/F,11/4/G,7/6/H}
		%		\tkzDefPoints{0/1/A,6/9/B, 3/0/C,9/0/D,11/3/E,16/3/F,16/8/G,12/10/H}
		\tkzDrawSegments(A,B  C,D D,E F,G G,H)
	\end{tikzpicture} 
	\caption{Quel chemin est le plus long ? (d'abord sans quadrillage, puis avec quadrillage et conclure en rajoutant le repère.) 
	}
\end{figure*} 

%----------------------------------------------------------------------------------------
%	 
%----------------------------------------------------------------------------------------
\newpage 
\section{Repères et repères orthonormés} 
%$\left( {{\mathrm{O}};\vec{\imath},\vec{\jmath}} \right)$ 
$O$, $I$ et $J$ trois points non alignés.
Un repère  $(O;I,J)$ ( ou $(O; \vv{OI}, \vv{OJ})$), du plan est formé de :
\begin{dingautolist}{192}
	\item origine $O$ du repère
	\item $(OI)$ est l'axe gradué des abscisse  (des $x$)
	\item $(OJ)$ est l'axe gradué des ordonnées $(OJ)$
\end{dingautolist}
Chaque point $M$ du plan correspond à un unique couple de coordonnées $\big(x(M); y(M)\big)\in\mathbb{R}\times\mathbb{R}$. Ces coordonnées se lisent sur les deux axes en tracant leurs parallèles passant par $M$.



\begin{figure*}[hb]\centering 
	\begin{tikzpicture}[scale=0.9] 
		\def \xmin{-1.5}; 
		\def \xmax{5.0}; 
		\def \ymin{-2};
		\def \ymax{5};	
		\def \alpha{0}
		\def \beta{90}
		\def \xnorm{1.0};
		\def \ynorm{1.0}; 
		\pgftransformcm{\xnorm*cos(\alpha)}{\xnorm*sin(\alpha)}{\ynorm*cos(\beta)}{\ynorm*sin(\beta)}{\pgfpointorigin}
		\path[clip,preaction = {draw=gray}] 
		({(\xmax*\ynorm*sin(\beta)-\ymax*\ynorm*cos(\beta))/(\xnorm*\ynorm*sin(\beta-\alpha))},{(-\xmax*\xnorm*sin(\alpha)+\ymax*\xnorm*cos(\alpha))/(\xnorm*\ynorm*sin(\beta-\alpha))}	) -- ( 
		{(\xmax*\ynorm*sin(\beta)-\ymin*\ynorm*cos(\beta))/(\xnorm*\ynorm*sin(\beta-\alpha))},{(-\xmax*\xnorm*sin(\alpha)+\ymin*\xnorm*cos(\alpha))/(\xnorm*\ynorm*sin(\beta-\alpha))}) -- (
		{(\xmin*\ynorm*sin(\beta)-\ymin*\ynorm*cos(\beta))/(\xnorm*\ynorm*sin(\beta-\alpha))},{(-\xmin*\xnorm*sin(\alpha)+\ymin*\xnorm*cos(\alpha))/(\xnorm*\ynorm*sin(\beta-\alpha))}) -- (
		{(\xmin*\ynorm*sin(\beta)-\ymax*\ynorm*cos(\beta))/(\xnorm*\ynorm*sin(\beta-\alpha))},{(-\xmin*\xnorm*sin(\alpha)+\ymax*\xnorm*cos(\alpha))/(\xnorm*\ynorm*sin(\beta-\alpha))}) --   cycle; 
		\tkzInit[xmin=-8, ymin=-8,  xmax=8.5,ymax=9.5]
		%		\begin{scope}[dashed] 
		%			\tkzGrid[ xstep=1, ystep=1](-11,-9)(10,15)
		%		\end{scope}  
		\tkzDrawX[left=5pt, label=, font=\scriptsize, line width=1.2pt ]  
		\tkzDrawY[below =5pt, label=, font=\scriptsize, line width=1.2pt] 
		\tkzRep[color  = Maroon,xlabel=,ylabel=, line width = 2pt,xnorm=1, ynorm=1]  
		\tkzDefPoint(0,0){O} 
		\tkzDefPoint(1,0){I} 
		\tkzDefPoint(0,1){J} 
		\tkzLabelPoints[below left,  font=\scriptsize](O)
		\tkzLabelPoints[below,  font=\scriptsize](I)
		\tkzLabelPoints[left,  font=\scriptsize](J)
		\tkzDefPoint(2.3,3.7){M}
		\tkzPointShowCoord[xlabel=$x$,ylabel=$y$](M)
		\tkzLabelPoint[above](M){\scriptsize$M(x;y)$} 
		\tkzMarkRightAngle[fill=lightgray, size=0.3](I,O,J);
		\tkzMarkSegments[mark=s||](O,I O,J)
	\end{tikzpicture}
	\begin{tikzpicture}[scale=0.9] 
		\def \xmin{-1.5}; 
		\def \xmax{5.0}; 
		\def \ymin{-2};
		\def \ymax{5};	
		\def \alpha{0}
		\def \beta{90}
		\def \xnorm{1.5};
		\def \ynorm{0.8}; 
		\pgftransformcm{\xnorm*cos(\alpha)}{\xnorm*sin(\alpha)}{\ynorm*cos(\beta)}{\ynorm*sin(\beta)}{\pgfpointorigin}
		\path[clip,preaction = {draw=gray}] 
		({(\xmax*\ynorm*sin(\beta)-\ymax*\ynorm*cos(\beta))/(\xnorm*\ynorm*sin(\beta-\alpha))},{(-\xmax*\xnorm*sin(\alpha)+\ymax*\xnorm*cos(\alpha))/(\xnorm*\ynorm*sin(\beta-\alpha))}	) -- ( 
		{(\xmax*\ynorm*sin(\beta)-\ymin*\ynorm*cos(\beta))/(\xnorm*\ynorm*sin(\beta-\alpha))},{(-\xmax*\xnorm*sin(\alpha)+\ymin*\xnorm*cos(\alpha))/(\xnorm*\ynorm*sin(\beta-\alpha))}) -- (
		{(\xmin*\ynorm*sin(\beta)-\ymin*\ynorm*cos(\beta))/(\xnorm*\ynorm*sin(\beta-\alpha))},{(-\xmin*\xnorm*sin(\alpha)+\ymin*\xnorm*cos(\alpha))/(\xnorm*\ynorm*sin(\beta-\alpha))}) -- (
		{(\xmin*\ynorm*sin(\beta)-\ymax*\ynorm*cos(\beta))/(\xnorm*\ynorm*sin(\beta-\alpha))},{(-\xmin*\xnorm*sin(\alpha)+\ymax*\xnorm*cos(\alpha))/(\xnorm*\ynorm*sin(\beta-\alpha))}) --   cycle; 
		\tkzInit[xmin=-8, ymin=-8,  xmax=8.5,ymax=9.5]
		%		\begin{scope}[dashed] 
		%			\tkzGrid[ xstep=1, ystep=1](-11,-9)(10,15)
		%		\end{scope}  
		\tkzDrawX[left=5pt, label=, font=\scriptsize, line width=1.2pt ]  
		\tkzDrawY[below =5pt, label=, font=\scriptsize, line width=1.2pt] 
		\tkzRep[color  = Maroon,xlabel=,ylabel=, line width = 2pt,xnorm=1, ynorm=1]  
		\tkzDefPoint(0,0){O} 
		\tkzDefPoint(1,0){I} 
		\tkzDefPoint(0,1){J} 
		\tkzLabelPoints[below left,  font=\scriptsize](O)
		\tkzLabelPoints[below,  font=\scriptsize](I)
		\tkzLabelPoints[left,  font=\scriptsize](J)
		\tkzDefPoint(2.3,3.7){M}
		\tkzPointShowCoord[xlabel=$x$,ylabel=$y$](M)
		\tkzLabelPoint[above](M){\scriptsize$M(x;y)$}
	\end{tikzpicture}
	\begin{tikzpicture}[scale=0.9] 
		\def \xmin{-1.5}; 
		\def \xmax{5.0}; 
		\def \ymin{-2};
		\def \ymax{5};	
		\def \alpha{-05}
		\def \beta{75}
		\def \xnorm{1.5};
		\def \ynorm{0.8}; 
		\pgftransformcm{\xnorm*cos(\alpha)}{\xnorm*sin(\alpha)}{\ynorm*cos(\beta)}{\ynorm*sin(\beta)}{\pgfpointorigin}
		\path[clip,preaction = {draw=gray}] 
		({(\xmax*\ynorm*sin(\beta)-\ymax*\ynorm*cos(\beta))/(\xnorm*\ynorm*sin(\beta-\alpha))},{(-\xmax*\xnorm*sin(\alpha)+\ymax*\xnorm*cos(\alpha))/(\xnorm*\ynorm*sin(\beta-\alpha))}	) -- ( 
		{(\xmax*\ynorm*sin(\beta)-\ymin*\ynorm*cos(\beta))/(\xnorm*\ynorm*sin(\beta-\alpha))},{(-\xmax*\xnorm*sin(\alpha)+\ymin*\xnorm*cos(\alpha))/(\xnorm*\ynorm*sin(\beta-\alpha))}) -- (
		{(\xmin*\ynorm*sin(\beta)-\ymin*\ynorm*cos(\beta))/(\xnorm*\ynorm*sin(\beta-\alpha))},{(-\xmin*\xnorm*sin(\alpha)+\ymin*\xnorm*cos(\alpha))/(\xnorm*\ynorm*sin(\beta-\alpha))}) -- (
		{(\xmin*\ynorm*sin(\beta)-\ymax*\ynorm*cos(\beta))/(\xnorm*\ynorm*sin(\beta-\alpha))},{(-\xmin*\xnorm*sin(\alpha)+\ymax*\xnorm*cos(\alpha))/(\xnorm*\ynorm*sin(\beta-\alpha))}) --   cycle; 
		\tkzInit[xmin=-8, ymin=-8,  xmax=8.5,ymax=9.5]
		%		\begin{scope}[dashed] 
		%			\tkzGrid[ xstep=1, ystep=1](-11,-9)(10,15)
		%		\end{scope}  
		\tkzDrawX[left=5pt, label=, font=\scriptsize, line width=1.2pt ]  
		\tkzDrawY[below =5pt, label=, font=\scriptsize, line width=1.2pt] 
		\tkzRep[color  = Maroon,xlabel=,ylabel=, line width = 2pt,xnorm=1, ynorm=1]  
		\tkzDefPoint(0,0){O} 
		\tkzDefPoint(1,0){I} 
		\tkzDefPoint(0,1){J} 
		\tkzLabelPoints[below left,  font=\scriptsize](O)
		\tkzLabelPoints[below,  font=\scriptsize](I)
		\tkzLabelPoints[left,  font=\scriptsize](J)
		\tkzDefPoint(2.3,3.7){M}
		\tkzPointShowCoord[xlabel=$x$,ylabel=$y$](M)
		\tkzLabelPoint[above](M){\scriptsize$M(x;y)$}
	\end{tikzpicture}
	\caption{Dans un repère orthonormé $(O;I,J)$, le triangle $OIJ$ est rectangle isocèle (gauche). Le repère $(O;I,J)$ est orthogonal si $OIJ$ est rectangle mais pas isocèle (milieu). Les repères peuvent être obliques (droite).}
\end{figure*}


Par la suite, on travaille avec des \textbf{repères orthonormés} : 
\begin{description}
	\item[ortho-gonal]  les axes sont perpendiculaires
	\item[normé] la longueur des segments unitaires sur les deux axes est la mêmes
\end{description}
%$(O;I,J)$ est orthonormé losrque le triangle $OIJ$ est rectangle isocèle en $O$.

%----------------------------------------------------------------------------------------
%	 
%----------------------------------------------------------------------------------------
\newpage 
\section{Distance dans les repères orthonormés}

\begin{figure}[hb] 
	\begin{tikzpicture}[scale=0.90]
		\tkzSetUpPoint[size=3, fill=black]
		\tkzfctset{tan style/.style={->,>=latex, color=red, line width=1.2pt}}   
		\tkzSetUpLine[line width= 1.2pt, add=.5 and .5] 
		
		\def \xmin{-2}; 
		\def \xmax{11}; 
		\def \ymin{-1.5};
		\def \ymax{9};	
		\def \alpha{0}
		\def \beta{90}
		\def \xnorm{1.0};
		\def \ynorm{1.0}; 
		\pgftransformcm{\xnorm*cos(\alpha)}{\xnorm*sin(\alpha)}{\ynorm*cos(\beta)}{\ynorm*sin(\beta)}{\pgfpointorigin}
		\path[clip,preaction = {draw=gray}] 
		({(\xmax*\ynorm*sin(\beta)-\ymax*\ynorm*cos(\beta))/(\xnorm*\ynorm*sin(\beta-\alpha))},{(-\xmax*\xnorm*sin(\alpha)+\ymax*\xnorm*cos(\alpha))/(\xnorm*\ynorm*sin(\beta-\alpha))}	) -- ( 
		{(\xmax*\ynorm*sin(\beta)-\ymin*\ynorm*cos(\beta))/(\xnorm*\ynorm*sin(\beta-\alpha))},{(-\xmax*\xnorm*sin(\alpha)+\ymin*\xnorm*cos(\alpha))/(\xnorm*\ynorm*sin(\beta-\alpha))}) -- (
		{(\xmin*\ynorm*sin(\beta)-\ymin*\ynorm*cos(\beta))/(\xnorm*\ynorm*sin(\beta-\alpha))},{(-\xmin*\xnorm*sin(\alpha)+\ymin*\xnorm*cos(\alpha))/(\xnorm*\ynorm*sin(\beta-\alpha))}) -- (
		{(\xmin*\ynorm*sin(\beta)-\ymax*\ynorm*cos(\beta))/(\xnorm*\ynorm*sin(\beta-\alpha))},{(-\xmin*\xnorm*sin(\alpha)+\ymax*\xnorm*cos(\alpha))/(\xnorm*\ynorm*sin(\beta-\alpha))}) --   cycle;  
		
		\tkzInit[xmin=-8,xmax=10,ymin=-7,ymax=8  ]
		\tkzDrawX[  noticks,label={$x$}, font=\scriptsize, line width=1.2pt ]  
		\tkzDrawY[ noticks, label={$y$}, font=\scriptsize, line width=1.2pt] 
		\tkzRep[color  = Maroon,xlabel=,ylabel=, font=\scriptsize, line width = 2pt,xnorm=1, ynorm=1]   
		%		\begin{scope}[dashed] 
		%			\tkzGrid[ xstep=1, ystep=1](-2,-1.5)(11,9)
		%		\end{scope}
		\tkzDefPoint(0,0){O} 
		\tkzDefPoint(1,0){I} 
		\tkzDefPoint(0,1){J}  
		\tkzDrawPoints(O,I,J)
		\tkzLabelPoints[font=\scriptsize](I)
		\tkzLabelPoints[left, font=\scriptsize](J)
		\tkzLabelPoints[below left, font=\scriptsize](O)  
		
		
		\tkzDefPoint(3,4){A};
		\tkzDefPoint(8,6){B};
		\tkzDefPoint(8, 4){C}; 
		\tkzDrawSegment[dim={\scriptsize$\abs{x_B-x_A}$,-15pt, ->, >=latex },dashed](A,C))
		\tkzDrawSegment[dim={\scriptsize$\abs{y_B-y_A}$,-25pt, ->, >=latex },dashed](C,B))
		\tkzDrawSegment(A,B)
		\tkzPointShowCoord[   xlabel={\color{blue}$x_A$}, xstyle={below} , ylabel={\color{red}$y_A$}, ystyle={left} ](A)
		\tkzPointShowCoord[   xlabel={\color{blue}$x_B$}, xstyle={below} , ylabel={\color{red}$y_B$}, ystyle={left} ](B) 
		\tkzMarkRightAngle[fill=lightgray, size=0.4](B,C,A);
		
		\tkzDrawPoints(A,B)
		\tkzLabelPoints[above, font=\scriptsize ](A)
		\tkzLabelPoints[above, font=\scriptsize](B)
	\end{tikzpicture}
	\caption{ 
		Le plan est muni d'un repère \textbf{orthonormé} $(O,I,J)$. \\
		$ABC$ est rectangle en $C$.  
	}
\end{figure}  
\begin{proof}
	Le repère est orthonormé, le triangle $ABC$ rectangle en $C$. D'après théorème de Pythagore : 
	\begin{align*}
		AC&=|x_{B}-x_{A}| & 
		BC&=|y_{B}-y_{A}|	\\
		AB^2&= AC^2 + CB^2 \\
		AB^2& = \left|x_{B}-x_{A}\right|^2+\left|y_{B}-y_{A}\right|^2\\
		AB^2& =\left(x_{B}-x_{A}\right)^2+\left(y_{B}-y_{A}\right)^2 &&\text{$\abs{a}^2=a^2$}\\
		AB & =\sqrt{\left(x_{B}-x_{A}\right)^2+\left(y_{B}-y_{A}\right)^2}
	\end{align*} 
\end{proof} 
\begin{theorem} Dans un repère orthonormé. La distance entre les points $A(x_A;y_A)$ et $B(x_B;y_B)$ est donnée par les formules :
	$$AB^2= \left(x_{B}-x_{A}\right)^2+\left(y_{B}-y_{A}\right)^2$$
	$$AB   =\sqrt{\left(x_{B}-x_{A}\right)^2+\left(y_{B}-y_{A}\right)^2}$$
\end{theorem} 
\newpage 
\begin{example}
	Dans un repère orthonormé, $A(1;5)$ et $B(-2;7)$ :
	
	
	\QCM[VF,Alterne,Solution,Largeur=1.0cm,Stretch=1.3]{%
		{$AB^2= (1+2)^2-(5-7)^2=$ }&2,%  
		{$AB^2= (1-2)^2+(5-7)^2=$ }&2,%
		{$AB^2= (1-5)^2+(-2-7)^2=$ }&2,% 
		{$AB^2= (1+2)^2+(5-7)^2=$ }&1,% 
		{$AB^2= (1-(-2))^2+(5-7)^2=$ }&1,% 
		{$AB^2= (-2-1)^2+(7-5)^2=$ }&1,%  
	} 
	$$
	AB=\sqrt{3^2+2^2}=\sqrt{13}
	$$
\end{example}
\begin{example}  Soit $M$ un point du cercle $\mathscr{C}(O;7)$ est le cercle de centre $O(0;0)$ et de rayon $7$. 
	\begin{align*}
		M(x;y)\in \mathscr{C}(O;7) & \iff OM = 7 \\
		& \iff OM^2 = 49\\
		& \iff \big( x - 0 \big)^2 + \big( y - 0 \big)^2 =49 \\
		& \iff x^2 + y^2 =49 
	\end{align*} 
	L'équation $x^2 + y^2 =49$ sur les coordonnées $(x;y)$ caractérise les points du cercle. 
	
	On dira que le cercle $\mathscr{C}$ a pour équation $x^2 + y^2 =49$. En abrégé :
	\[
	\mathscr{C}(O;7)\colon x^2 + y^2 =49
	\] 
\end{example}
